\documentclass[12pt]{article}

% Core math
\usepackage{amsmath,amssymb,amsfonts}
\usepackage{amsthm}

% Diagrams
\usepackage{tikz-cd}
\usepackage{tikz}
\usetikzlibrary{calc}

% Tables
\usepackage{array}
\usepackage{booktabs}

% Spacing helpers
\usepackage{needspace}
\usepackage{placeins}   % \FloatBarrier to keep Table 1 with its section

% References
\usepackage{hyperref}
\usepackage{cleveref}

% Graphics
\usepackage{graphicx}

% Page layout
\usepackage[margin=1in]{geometry}

% GrokRxiv sidebar
\usepackage{xcolor}

% ==== topological-phases-of-matter : canonical macros (v1, 2026-07-14) ====
% Requires amsmath, amssymb, amsfonts (mathfrak, mathbf, mathcal all standard).

% --- condensation / probes ---
\newcommand{\cond}[1]{\underline{#1}}          % condensation X |-> \cond{X}
\newcommand{\cg}[1]{\underline{#1}}            % condensed group (alias, for readability)
\newcommand{\Cont}{\operatorname{Cont}}        % Cont(S,X)
\newcommand{\Shape}{\operatorname{Shape}}      % shape / homotopy type

% --- interaction space & locality ---
\newcommand{\Int}{\mathcal{I}}                 % \Int_{d,G} interaction space
\newcommand{\Ffun}{F}                          % F-function (interaction decay)
\newcommand{\BF}{\mathcal{B}_{F}}              % interaction Banach space for F

% --- moduli stacks ---
\newcommand{\Ham}{\mathfrak{Ham}}              % \Ham_{d,G}
\newcommand{\Gap}{\mathfrak{Gap}}              % \Gap_{d,G}
\newcommand{\Phase}{\mathfrak{Phase}}          % \Phase_{d,G}
\newcommand{\Phases}{\operatorname{Phases}}    % \Phases_{d,G} = pi_0 Shape Phase
\newcommand{\Weq}{\mathcal{W}}                 % class of equivalences to invert
\newcommand{\st}{\mathrm{st}}                  % stabilization superscript
\newcommand{\hfix}[1]{h#1}                     % homotopy-fixed-point superscript, e.g. \Phase_d^{\hfix{\cg{G}}}
\newcommand{\quotstack}[2]{[#1/#2]}            % action stack [Y/\cg{G}]

% --- gap ---
\newcommand{\gapof}[1]{\operatorname{gap}(#1)} % gap(H)
\newcommand{\gapbound}{\Delta}                  % uniform gap bound

% --- stacking / spectrum ---
% \boxtimes is built in (stacking product)
\newcommand{\IPcond}{\mathbf{IP}^{\mathrm{cond}}}  % invertible condensed phase spectrum

% --- transitions / discriminant / charge ---
\newcommand{\Sig}{\Sigma}                       % \Sig_f gapless discriminant
\newcommand{\rc}{\partial\nu}                   % relative transition charge in E^{q+1}(B,U_f)

% --- observables & solid K-theory ---
\newcommand{\Kop}{K_{\mathrm{op}}}             % operator / topological K-theory
\newcommand{\Kalg}{K_{\mathrm{alg}}}           % algebraic K-theory
\newcommand{\Solid}{\operatorname{Solid}}      % solidification functor

% --- disorder ---
\newcommand{\dis}{\Omega}                       % disorder hull \Omega = F^{Z^d}
% ==== end canonical macros ====

% --- local convenience macros (not overriding canonical glyphs) ---
\newcommand{\RR}{\mathbb{R}}
\newcommand{\CC}{\mathbb{C}}
\newcommand{\ZZ}{\mathbb{Z}}
\newcommand{\NN}{\mathbb{N}}
\newcommand{\Uone}{U(1)}
\newcommand{\Hilb}{\mathcal{H}}
\newcommand{\qA}{\mathcal{A}}                  % quasi-local algebra
\newcommand{\St}{\mathcal{S}}                  % state space S(A)
\newcommand{\Aut}{\operatorname{Aut}}
\newcommand{\SRE}{\mathrm{SRE}}
\newcommand{\Groth}{\operatorname{Groth}}      % group completion
\newcommand{\IFH}{I_{\mathrm{FH}}}             % Freed--Hopkins invertible-TQFT spectrum
\newcommand{\real}{\operatorname{real}}        % realization/comparison map
\newcommand{\CPr}{\mathrm{cp}}                  % commuting-projector superscript
\newcommand{\Obs}{\operatorname{Obs}}
\newcommand{\norm}[1]{\lVert #1 \rVert}
\newcommand{\abs}[1]{\lvert #1 \rvert}
\newcommand{\set}[1]{\{\,#1\,\}}
\newcommand{\dist}{\operatorname{dist}}
\newcommand{\im}{\operatorname{im}}
\newcommand{\id}{\mathrm{id}}
\DeclareMathOperator*{\einf}{inf}

% Theorem environments. We use aliascnt so that each environment carries its own
% counter (aliased to the Theorem counter for shared numbering); this lets cleveref
% name cross-references correctly ("Proposition 2.3", "Recollection 3.1") instead of
% collapsing them all to "Theorem".
\usepackage{aliascnt}
\theoremstyle{plain}
\newtheorem{theorem}{Theorem}[section]
\newaliascnt{proposition}{theorem}
\newtheorem{proposition}[proposition]{Proposition}
\aliascntresetthe{proposition}
\newaliascnt{lemma}{theorem}
\newtheorem{lemma}[lemma]{Lemma}
\aliascntresetthe{lemma}
\newaliascnt{corollary}{theorem}
\newtheorem{corollary}[corollary]{Corollary}
\aliascntresetthe{corollary}
\newaliascnt{recollection}{theorem}
\newtheorem{recollection}[recollection]{Recollection}
\aliascntresetthe{recollection}
% Numbered program-level conjectures, tagged VI-k per the series convention
% (own counter so cross-references resolve to "Conjecture VI-k", not the section).
\newtheorem{progconj}{Conjecture}
\renewcommand{\theprogconj}{VI-\arabic{progconj}}
\theoremstyle{definition}
\newaliascnt{definition}{theorem}
\newtheorem{definition}[definition]{Definition}
\aliascntresetthe{definition}
\newaliascnt{example}{theorem}
\newtheorem{example}[example]{Example}
\aliascntresetthe{example}
\theoremstyle{remark}
\newaliascnt{remark}{theorem}
\newtheorem{remark}[remark]{Remark}
\aliascntresetthe{remark}
% cleveref display names for the aliased environments.
\crefname{theorem}{Theorem}{Theorems}
\Crefname{theorem}{Theorem}{Theorems}
\crefname{proposition}{Proposition}{Propositions}
\Crefname{proposition}{Proposition}{Propositions}
\crefname{lemma}{Lemma}{Lemmas}
\Crefname{lemma}{Lemma}{Lemmas}
\crefname{corollary}{Corollary}{Corollaries}
\Crefname{corollary}{Corollary}{Corollaries}
\crefname{recollection}{Recollection}{Recollections}
\Crefname{recollection}{Recollection}{Recollections}
\crefname{definition}{Definition}{Definitions}
\Crefname{definition}{Definition}{Definitions}
\crefname{example}{Example}{Examples}
\Crefname{example}{Example}{Examples}
\crefname{remark}{Remark}{Remarks}
\Crefname{remark}{Remark}{Remarks}
\crefname{progconj}{Conjecture}{Conjectures}
\Crefname{progconj}{Conjecture}{Conjectures}

% GrokRxiv DOI sidebar
\definecolor{grokgray}{RGB}{110,110,110}
\AddToHook{shipout/foreground}{%
  \ifnum\value{page}=1
    \begin{tikzpicture}[remember picture, overlay]
      \node[
        rotate=90,
        anchor=south,
        font=\Large\sffamily\bfseries\color{grokgray},
        inner sep=0pt
      ] at ([xshift=38pt, yshift=0.52\paperheight]current page.south west)
      {GrokRxiv:2026.07.synthesis\quad
       [\,math-ph\,]\quad
       14 Jul 2026};
    \end{tikzpicture}
  \fi
}

\hypersetup{
  colorlinks=true,
  linkcolor=blue!55!black,
  citecolor=green!45!black,
  urlcolor=blue!55!black
}

\title{\textbf{Topological Phases of Matter in the Condensed-Mathematics Paradigm:\\
A Modular Research Program}\\[0.4em]
\large Part VI of \emph{Topological Phases of Matter in the Condensed-Mathematics Paradigm}}

\author{Matthew Long \\
\textit{The YonedaAI Collaboration} \\
\textit{YonedaAI Research Collective} \\
Chicago, IL \\
\texttt{matthew@yonedaai.com} $\cdot$ \url{https://yonedaai.com}}

\date{14 July 2026}

\begin{document}
\maketitle

\begin{abstract}
This synthesis assembles the five analytic and topological modules of the preceding
Parts into one research program: topological phases of matter, read as the components
of a condensed moduli stack of uniformly gapped local Hamiltonians. The spine is a
single sequence:
local quantum interactions $\rightsquigarrow$ condensed moduli stack $\Ham_{d,G}$
$\rightsquigarrow$ uniformly gapped substack $\Gap_{d,G}$ $\rightsquigarrow$ stabilized
phase $\infty$-groupoid $\Phase_{d,G}$ $\rightsquigarrow$ invertible condensed phase
spectrum $\IPcond_{d,G}$, and each arrow is owned by a Part. Part~I makes the
interaction space a condensed Banach object and the Heisenberg flow a condensed
morphism with a base-uniform Lieb--Robinson light cone; Part~II attaches the
quasi-local $C^\ast$-algebra, its compact condensed state space, and the solid
$K$-theory invariant obtained from Aoki's solidification theorem (correctly, only
after inverting the Bott class); Part~III cuts out the uniformly gapped substack and
proves its stability on the frustration-free / LTQO stratum against the wall of
Cubitt--P\'erez-Garc\'ia--Wolf undecidability; Part~IV group-completes the invertible
sector into the condensed spectrum and sets it against the Freed--Hopkins bordism
target; Part~V asks which spectral classes a lattice model realizes, bounded by the
commuting-projector no-gos of Kapustin--Fidkowski and Kapustin--Spodyneiko. We treat the
program as a \emph{modular composition}: each
Part is a module whose output is the next module's input, and structure emerges at
each level: functorial invariants of disordered families, a robust set of phases, the
higher homotopy of pumps and defects, an obstruction subgroup of unrealized classes.
We develop the transition calculus, in which a topological phase transition is the
crossing of the gapless discriminant $\Sig_f$ and its charge is the relative class
$\rc\in E^{q+1}(B,B\setminus\Sig_f)$, an obstruction detected by linking spheres; we
work the SSH chain through all five modules; and we run the profinite disorder thread
$\dis=\Ffun^{\ZZ^d}$ through the whole. The module-level results are theorems, cited to
Parts~I--V\@. The program-level statements (that the modules glue into a single
condensed higher stack with $\IPcond_{d,G}$ its invertible spectrum, that the
microscopic and field-theoretic classifications agree, that the relative charge is a
spectral boundary map, that all invariants descend along disorder hulls) are stated as
numbered Conjectures~VI-1 through VI-4 and never as proven. A consolidated table
records all nineteen module conjectures with their dependencies. The genuinely new
thesis is not a different Chern number; it is a single categorical and homological
environment for continuous families, profinite disorder, completions, symmetry,
stacking, defects, and transition loci.
\end{abstract}

\tableofcontents

\section{Introduction}
\label{sec:intro}

\subsection{The proposal}
\label{sec:proposal}

A quantum phase is usually presented as an entry in a classification table. This
program takes a different starting point, due in outline to the seed prospectus of the
series and realized in detail across Parts~I--V\@: a phase is a \emph{component of a
condensed moduli object of uniformly gapped local Hamiltonians}. Fix a spatial
dimension $d$, a lattice or coarse metric space $L$, on-site Hilbert spaces, a locality
(interaction-decay) class, and an internal or spatial symmetry group $G$. Let
$\Int_{d,G}$ be the space of admissible $G$-symmetric local or quasi-local
interactions. Rather than use $\Int_{d,G}$ as an ordinary topological space, pass to
its condensation
\begin{equation}
\label{eq:condensation}
X\longmapsto\cond{X},\qquad \cond{X}(S)=\Cont(S,X),\qquad S\ \text{profinite},
\end{equation}
a sheaf on the site of profinite sets with finite jointly surjective covers. On
compactly generated spaces the passage $X\mapsto\cond{X}$ is fully faithful
\cite{clausen-scholze-condensed,barwick-haine-pyknotic}, so coupling tori, Banach
interaction spaces, and compact disorder hulls are not discarded but embedded into a
category with exact homological algebra. Because Hamiltonians carry gauge and
quasi-local automorphisms, the right object is a condensed \emph{anima} or
\emph{stack}, not a condensed set.

The whole program is the study of one geometric object and one substack. Over a
profinite probe $S$ let $\Ham_{d,G}(S)$ be the $S$-continuous families of
$G$-symmetric quasi-local interactions, and let
\begin{equation}
\label{eq:gapdef}
\Gap_{d,G}(S)=\bigcup_{\gapbound>0}\set{H\in\Ham_{d,G}(S):\einf_{s\in S}\gapof{H_s}\ge\gapbound}
\end{equation}
be the \emph{uniformly} gapped substack (the word uniformly being load-bearing, since
a family whose pointwise gaps degenerate to zero must be excluded). Inverting a class
$\Weq$ of gapped adiabatic / quasi-local equivalences and stabilizing by trivial
ancillas produces the phase $\infty$-groupoid, and its shape records the phases:
\begin{equation}
\label{eq:phase}
\Phase_{d,G}:=\bigl(\Gap_{d,G}[\Weq^{-1}]\bigr)^{\st},
\qquad
\Phases_{d,G}=\pi_0\,\Shape(\Phase_{d,G}).
\end{equation}
Stacking systems, $\boxtimes$, makes $\Phase_{d,G}$ symmetric monoidal; group-completing
its invertible sector yields the connective condensed spectrum $\IPcond_{d,G}$. The two
slogans of the program are the anchors of everything below:
\begin{equation}
\label{eq:slogans}
\boxed{\;\text{topological phase}=\text{a component of the stabilized condensed stack of gapped systems}\;}
\end{equation}
and
\begin{equation}
\label{eq:slogan2}
\boxed{\;
\begin{array}{c}
\text{topological transition}={}\\[2pt]
\text{crossing the gapless discriminant in the Hamiltonian moduli stack}.
\end{array}
\;}
\end{equation}

\subsection{The program as a modular hierarchy}
\label{sec:modular}

What Parts~I--V supply is not a monolith but five interoperable \emph{modules}, each of
which takes the previous module's output as its input and produces new structure. This
is the sense in which the program is modular: the composition is a hierarchy, and at
each level of the hierarchy a property emerges that was not visible one level down.

\begin{enumerate}
\item \textbf{Locality (Part I \cite{paperLocality}).} Makes $\Int_{d,G}$ a real Banach
  space $\BF$ through a Nachtergaele--Sims--Young $\Ffun$-function locality norm, and
  promotes the Heisenberg dynamics to a morphism of condensed sets with a
  \emph{base-uniform} Lieb--Robinson light cone. Emergent structure: the ground floor
  $\Ham_{d,G}$ exists as a condensed object with a well-defined notion of quasi-local
  time evolution over any profinite base.
\item \textbf{Positivity and $C^\ast$-norms (Part II \cite{paperPositivity}).} Attaches
  the quasi-local $C^\ast$-algebra $\qA$, its compact condensed state space, and, via
  Aoki's solidification theorem, a \emph{solid} $K$-theory invariant. Emergent
  structure: functorial topological invariants of disordered families over profinite
  hulls, produced by an intrinsically condensed operation.
\item \textbf{The gap (Part III \cite{paperGap}).} Cuts out the uniformly gapped
  substack $\Gap_{d,G}$ and proves its stability where a stability theory holds with
  uniform constants. Emergent structure: a robust set of phases $\Phases_{d,G}$, once
  $\Weq$ is inverted and ancillas are added.
\item \textbf{Stabilization and effective field theory (Part IV \cite{paperEFT}).}
  Group-completes the invertible sector into the connective condensed spectrum
  $\IPcond_{d,G}$ and compares it with the Freed--Hopkins bordism classification.
  Emergent structure: the higher homotopy of pumps and defects, and the relative charge
  $\rc$ as a spectral boundary map.
\item \textbf{Realizability (Part V \cite{paperRealizability}).} Identifies which
  spectral classes are produced by an explicit gapped lattice Hamiltonian. Emergent
  structure: an obstruction subgroup separating abstract classes from lattice-realizable
  ones.
\end{enumerate}

Two threads cut across all five modules and are what make the composition cohere rather
than merely stack. The first is the \emph{transition calculus}
$(U_f,\Sig_f,\nu_f,\rc)$, which ties modules III--V together: the gapped locus of a
family, its gapless discriminant, the locally constant phase label, and the relative
charge of a crossing. The second is the \emph{profinite disorder thread}
$\dis=\Ffun^{\ZZ^d}$, which runs through every module: descent of Lieb--Robinson
estimates in Part~I, crossed-product observable algebras in Part~II, uniform-gap descent
in Part~III, parametrized invariants in Part~IV, and functorial realizability in
Part~V\@.

\subsection{What is proved and what is conjectured}
\label{sec:stance}

The stance of this synthesis is the stance of the whole series: a rigorous research
\emph{program}, not a finished theory. We are strict about the boundary. The
\emph{module-level} results (Theorems~I-A through V-C, and the citations of Aoki,
Kubota, Freed--Hopkins, Ogata, and Kapustin--Fidkowski as rigorous anchors) are
theorems, and we restate them (as Recollections, with citations to the Parts that prove
them) rather than reprove them. The \emph{program-level} claims are open. That the
modules glue into a single condensed higher stack with $\IPcond_{d,G}$ as its invertible
spectrum; that the microscopic and field-theoretic classifications agree; that the
relative charge is the boundary map of that spectrum; that every program invariant
descends along disorder hulls; each of these is stated below as a numbered
Conjecture~VI-$k$ and is never asserted as proven.

Two hard theorems bound the program from outside and are quoted wherever they bite.
Cubitt--P\'erez-Garc\'ia--Wolf \cite{cpw-undecidable} proved that the spectral gap is
undecidable, so there is no algorithm and no uniform criterion deciding membership in
$\Gap_{d,G}$; existence of a thermodynamic gap is a \emph{hypothesis that defines} the
substack, not a computable property. Two commuting-projector no-gos bound realizability:
Kapustin--Fidkowski \cite{kapustin-fidkowski} proved that a $U(1)$-symmetric local
commuting-projector Hamiltonian has zero electric Hall conductance ($\sigma_H=0$), and
Kapustin--Spodyneiko \cite{kapustin-spodyneiko-thermal} proved that its chiral central charge
vanishes ($c_-=0$) with no symmetry hypothesis; so every chiral invertible class (nonzero
$\sigma_H$ under $U(1)$, or nonzero $c_-$) lies outside the commuting-projector image.
Neither wall is ever crossed in what follows.

\subsection{Relation to companion papers}
\label{sec:companions}

This paper is the capstone of a six-part series and consumes the outputs of Parts~I--V\@;
it imports the canonical notation of the series verbatim and does not restate it. We
summarize each companion in one paragraph, and \Cref{sec:spine} develops the
composition in full.

\emph{Part~I, Condensed Locality} \cite{paperLocality}, is the analytic ground floor.
Fixing an $\Ffun$-function it makes the $G$-symmetric interactions of finite
$\Ffun$-norm a real Banach space $\BF$, condenses it to $\cond{\BF}(S)=\Cont(S,\BF)$,
and proves that Heisenberg dynamics is a strongly continuous group of $\ast$-automorphisms
obeying a Lieb--Robinson bound (Theorem~I-A), that its assignment is a morphism of light
condensed sets (Theorem~I-B), and that a profinite family is exactly compatible finite
data with a base-uniform light cone (Theorem~I-C)\@. It generates the equivalence class
$\Weq$ from quasi-adiabatic continuation and states Conjectures~I-1 (condensed higher
stack), I-2 (condensed group of quasi-local automorphisms), and I-3 (descent of
Lieb--Robinson estimates along a disorder hull).

\emph{Part~II, Positivity, $C^\ast$-Norms, and Condensed State Spaces}
\cite{paperPositivity}, is the observable side. It shows the state space
$\St(\qA)$ condenses to a compact Hausdorff condensed set cut out by positivity and
normalization (Theorem~II-A), that Aoki's solidification recovers operator $K$-theory
after Bott inversion (Theorem~II-B), and that Bellissard's crossed product
$C(\dis)\rtimes\ZZ^d$ condenses functorially in the finite quotients of the hull
(Theorem~II-C)\@. It records Conjectures~II-1 through II-5 (closed condensed substack;
naturality of the solid invariant; Real/$KKO$ refinement; solid classifying object;
condensed superselection). Following Part~II we write the disorder alphabet as $Q$, not
$\Ffun$, to avoid collision with the $\Ffun$-function of Part~I; see
\Cref{sec:disorder}.

\emph{Part~III, The Uniformly Gapped Substack} \cite{paperGap}, is the gap module. It
proves that the finite-volume gap is Lipschitz and the gapped locus is open at finite
resolution (Theorem~III-A), that $\Gap_{d,G}$ is condensed-open on the frustration-free
/ LTQO stratum (Theorem~III-B), and that a uniform gap forces uniform clustering
(Theorem~III-C), all against the CPW undecidability wall; it states Conjectures~III-1
(openness on the physical
stratum), III-2 (uniform-gap descent), and III-3 (quasi-adiabatic continuation realizes
$\Weq$).

\emph{Part~IV, From Lattice Models to Effective Field Theories} \cite{paperEFT}, is the
stabilization module. It makes $\boxtimes$ a symmetric-monoidal structure, group-completes
$\Phases_{d,G}$, assembles $\IPcond_{d,G}$ with Kubota's $\Omega$-spectrum as its
rigorous carrier, and sets the lattice classification against Freed--Hopkins. It states
Conjectures~IV-1 (condensed/solid refinement), IV-2 (lattice--EFT comparison
equivalence), IV-3 (relative charge as spectral boundary map), IV-4 (solidification
commutes with group completion), and IV-5 (renormalization functor).

\emph{Part~V, Physical Realizability} \cite{paperRealizability}, is the realizability
module. It proves that every group-cohomology class is realized by a commuting-projector
model (Theorem~V-A), completeness of the operator-algebraic index in $d=1$ for on-site
finite symmetry, with the $d=2$ index surjective but its completeness open
(Theorem~V-B and Remark~V-B$'$), and the two commuting-projector no-gos, electric
(Kapustin--Fidkowski) and thermal/chiral (Kapustin--Spodyneiko) (Theorem~V-C)\@. It states
Conjectures~V-1 (realizability image is the SRE subspectrum),
V-2 (chiral realizability by non-commuting models), and V-3 (profinite-family
realizability).

\subsection{Outline}
\label{sec:outline}

\Cref{sec:substrate} recalls the condensed-mathematics substrate and fixes conventions.
\Cref{sec:spine} is the spine: the boxed sequence, arrow by arrow, with each module's
theorems restated and cited and its emergent structure named.
\Cref{sec:composition} makes modular composition explicit: the composition diagram, the
emergent properties, and the two walls. \Cref{sec:dictionary} states the homotopy
dictionary. \Cref{sec:transitions} develops the relative-charge formalism.
\Cref{sec:ssh} works the SSH chain through all five modules.
\Cref{sec:disorder} runs the profinite disorder thread. \Cref{sec:conjectures} states
the program-level Conjectures~VI-1 through VI-4 and consolidates all nineteen module
conjectures with dependencies. \Cref{sec:contributions} weighs what the paradigm
contributes and what it does not, and \Cref{sec:conclusion} concludes.

\section{The condensed-mathematics substrate}
\label{sec:substrate}

We recall only what the composition uses; Parts~I and II give the careful development.

\subsection{Condensation and profinite probes}

A condensed set is a sheaf on the site of profinite sets $S$ with covers the finite
jointly surjective families \cite{clausen-scholze-condensed}; the pyknotic variant of
Barwick--Haine \cite{barwick-haine-pyknotic} differs only in set-theoretic
bookkeeping. The functor \eqref{eq:condensation} embeds compactly generated topological
spaces fully faithfully, so no information is lost in passing from $X$ to $\cond{X}$.
The point of the passage is algebraic: condensed abelian groups form a Grothendieck
abelian category with exact products and derived functors, where ordinary topological
abelian groups do not, and this is precisely the environment in which derived
invariants, descent, and completions behave. A condensed object is \emph{light} when it
is generated under the site by a countable family (the case for separable Banach spaces
and separable $C^\ast$-algebras), which is why Parts~I and II take care to work with
separable data.

\subsection{Stacks, anima, and shape}

Because a Hamiltonian carries automorphisms (gauge transformations, on-site symmetry
actions, quasi-local unitaries), the moduli problem $S\mapsto\{S\text{-families of
Hamiltonians}\}$ valued in groupoids (or $\infty$-groupoids) is a candidate
\emph{condensed stack}, not merely a condensed set. Whether it satisfies descent along
profinite covers is the content of Conjecture~I-1 and is not proved. The \emph{shape}
$\Shape(-)$ of a condensed anima is its underlying homotopy type; $\pi_0\Shape$
extracts the ordinary set of components, hence $\Phases_{d,G}$ in \eqref{eq:phase}.
Throughout, fraktur denotes the moduli stacks ($\Ham,\Gap,\Phase$), calligraphic the
interaction space and the equivalence class ($\Int,\Weq$), boldface the spectra
($\IPcond$), and $\cond{(-)}$ condensation.

\subsection{The boxed sequence}

The program is summarized by one sequence, quoted from the prospectus and realized in
Parts~I--V\@:
\begin{equation}
\label{eq:boxseq}
\boxed{
\begin{array}{c}
\text{local quantum interactions}\\[2pt]
\big\downarrow\ \text{\footnotesize Part I}\\[2pt]
\text{condensed moduli stack of Hamiltonians }\Ham_{d,G}\\[2pt]
\big\downarrow\ \text{\footnotesize Part II}\\[2pt]
\text{condensed observable algebras and solid }K\text{-theory}\\[2pt]
\big\downarrow\ \text{\footnotesize Part III}\\[2pt]
\text{uniformly gapped substack }\Gap_{d,G}\\[2pt]
\big\downarrow\ \text{\footnotesize Part IV}\\[2pt]
\text{stabilized phase }\infty\text{-groupoid }\Phase_{d,G}\\[2pt]
\big\downarrow\ \text{\footnotesize Part IV}\\[2pt]
\text{invertible condensed phase spectrum }\IPcond_{d,G}\\[2pt]
\big\downarrow\ \text{\footnotesize Part V}\\[2pt]
\text{realized subspectrum inside }\IPcond_{d,G}.
\end{array}}
\end{equation}
The prospectus states the middle five lines; Parts~II and V add the observable-algebra
and realization refinements. Each arrow is a module. \Cref{sec:spine} treats them in
turn.

\section{The spine: the boxed sequence and its five modules}
\label{sec:spine}

We walk down \eqref{eq:boxseq}. For each module we recall its principal theorems (as
Recollections, attributed to the Part that proves them) and name the structure that
emerges. Nothing in this section is claimed as new; the novelty of the synthesis is the
composition, treated in \Cref{sec:composition}.

\subsection{Arrow I: local interactions to the condensed stack}
\label{sec:arrowI}

The first arrow turns a physicist's list of local couplings into a condensed object
that can be probed by profinite sets. Part~I fixes an $\Ffun$-function $\Ffun$ encoding
interaction decay and defines the interaction Banach space $\BF$.

\begin{recollection}[Banach interaction space and condensed dynamics; Part~I, Theorems~I-A and I-B \cite{paperLocality}]
\label{rec:I}
$\BF$ is a real Banach space; for $\Phi\in\BF$ the infinite-volume Heisenberg dynamics
$\tau^\Phi_t=\lim_{\Lambda\uparrow L}\tau^{\Phi,\Lambda}_t$ exists in operator norm,
uniformly for $t$ in compact sets, and is a strongly continuous one-parameter group of
$\ast$-automorphisms of the quasi-local algebra $\qA$ obeying a Lieb--Robinson bound. Its
condensation
\[
\cond{D}:\cond{\RR}\times\cond{\BF}\longrightarrow\cond{\Aut(\qA)},
\qquad (t,\Phi)\mapsto\tau^\Phi_t,
\]
is a morphism of (light, when the interactions are separable) condensed sets.
\end{recollection}

\begin{recollection}[Base-uniform light cone; Part~I, Theorem~I-C \cite{paperLocality}]
\label{rec:Ic}
Let $S=\varprojlim_i S_i$ be profinite and $\Phi_\bullet\in\Cont(S,\BF)=\cond{\BF}(S)$.
Then $\Phi_\bullet$ is a norm limit of families factoring through finite quotients
$S\to S_i$; $B=\sup_{s\in S}\norm{\Phi_s}_\Ffun<\infty$; and if
$\sup_{s}\norm{\Phi_s}_{\Ffun_a}<\infty$ for a reweighting $\Ffun_a(r)=e^{-ar}\Ffun(r)$,
the Lieb--Robinson bound holds for all $s\in S$ simultaneously with one velocity
$v=2\sup_s\norm{\Phi_s}_{\Ffun_a}C_{\Ffun_a}/a$.
\end{recollection}

\emph{Emergent structure.} The ground floor $\Ham_{d,G}$ is a condensed object over
which time evolution is quasi-local \emph{uniformly in the probe}. Uniformity is what
makes every later module's ``family over a base'' well-posed: it is the reason a
disordered family has a single light cone, a single clustering length, a single set of
constants. Part~I's Conjectures~I-1, I-2, I-3 mark the edge: stack descent, a condensed
automorphism group, and Lieb--Robinson descent along a hull; none proved.

\subsection{Arrow II: observable algebras and solid \texorpdfstring{$K$}{K}-theory}
\label{sec:arrowII}

The second arrow attaches to the stack the algebra of observables and the $K$-theory
from which topological invariants are built. For a spin system with finite on-site
dimension $\qA$ is a separable unital AF algebra (UHF in the homogeneous case), and its
condensation is a sheaf of $C^\ast$-algebras, light because $\qA$ is separable.

\begin{recollection}[Condensed state space; Part~II, Theorem~II-A \cite{paperPositivity}]
\label{rec:IIa}
For a unital $C^\ast$-algebra $A$, the state space $\cond{\St(A)}$ is a compact
Hausdorff condensed set on which condensation is fully faithful, and positivity together
with normalization exhibit it as a closed condensed subobject of the condensed dual
ball. On a uniformly gapped family the ground-state section is weak-$\ast$ continuous by
the Bachmann--Michalakis--Nachtergaele--Sims spectral-flow cocycle
\cite{bmns-automorphic}, a genuine point of $\cond{\St(A)}$ over the gapped locus.
\end{recollection}

The load-bearing bridge is Aoki's theorem, and the honesty of the whole program depends
on stating it with its Bott-inversion caveat, exactly as Part~II now does.

\begin{recollection}[Solidification and operator $K$-theory; Part~II, Theorem~II-B, after Aoki \cite{aoki-solidification,paperPositivity}]
\label{rec:aoki}
Let $A$ be a real associative algebra and $\cond{A}$ its condensation. Solidification of
the \emph{connective} algebraic $K$-theory of $\cond{A}$ is discrete and recovers the
connective semitopological $K$-theory of $A$ (Friedlander--Walker, and Blanc),
\[
\Solid\bigl(\Kalg(\cond{A})\bigr)\;\simeq\;K^{\mathrm{semi}}(A).
\]
If $A$ is a real Banach algebra this is the connective part of operator $K$-theory; the
full, Bott-periodic operator $K$-theory is recovered only after inverting the Bott class
$\beta$,
\begin{equation}
\label{eq:bottinv}
\Kop(A)\;\simeq\;\Solid\bigl(\Kalg(\cond{A})\bigr)\bigl[\beta^{-1}\bigr],
\end{equation}
with $\beta\in K_2^{\mathrm{top}}(\CC)\cong\ZZ$ complex, the period-eight real Bott class
in the $KO$ case (an infinite-order generator of the eightfold $KO$ periodicity).
\end{recollection}

We stress the bookkeeping because it is easy to get wrong and consequential when
composed. Writing $\Kop(A)\simeq\Solid(\Kalg(\cond{A}))$ without the inversion
$[\beta^{-1}]$ is false in negative degrees: solidification of connective algebraic
$K$-theory lands in a connective spectrum with no negative homotopy, whereas operator
$K$-theory is periodic. The two agree in nonnegative degrees, which is all the SSH
winding (a degree-$\ge 0$ class) needs, but the periodic identification requires
\eqref{eq:bottinv}. Every later use of ``the solid invariant'' in this synthesis carries
this qualification.

\begin{recollection}[Crossed-product disorder algebra; Part~II, Theorem~II-C \cite{paperPositivity}]
\label{rec:crossed}
For $\dis=Q^{\ZZ^d}$ with the shift action $T$, the covariant observable algebra of a
homogeneous disordered family is the crossed product $C(\dis)\rtimes_T\ZZ^d$
\cite{bellissard-ncg-qhe}, a separable unital $C^\ast$-algebra whose $K_0$-trace range is
the gap-labelling group; its condensation is functorial in the finite quotients of
$\dis$.
\end{recollection}

\emph{Emergent structure.} The composition of Arrows~I and II produces \emph{functorial
topological invariants of disordered families over profinite hulls}, computed by an
intrinsically condensed operation (solidification) rather than imported by hand. This is
the first place the condensed language does work an ordinary smooth-parameter treatment
does not: the profinite probes of \eqref{eq:condensation} match the physical
configuration space $\dis$ of \Cref{rec:crossed}, they do not approximate it. Part~II's
Conjectures~II-1 through II-5 record what is open here.

\subsection{Arrow III: the uniformly gapped substack}
\label{sec:arrowIII}

The third arrow selects the systems on which a phase invariant is even defined. The
selection is governed by a wall.

\begin{recollection}[Undecidability of the gap; Cubitt--P\'erez-Garc\'ia--Wolf \cite{cpw-undecidable}]
\label{rec:cpw}
There is a family of translation-invariant nearest-neighbour Hamiltonians on a
two-dimensional lattice, depending computably on a parameter, for which no algorithm
decides whether the thermodynamic-limit system is gapped or gapless; the property is
$\Pi_1$-complete.
\end{recollection}

Consequently membership in $\Gap_{d,G}$ is not, in general, decidable, and Part~III does
not attempt to decide it. Existence of a thermodynamic gap \emph{defines} the substack;
the content is stability.

\begin{recollection}[Gap continuity and stability; Part~III, Theorems~III-A and III-B \cite{paperGap}]
\label{rec:III}
At fixed finite volume the $n$-gap is $2C_\Lambda$-Lipschitz in the interaction norm, so
the finite-volume gapped locus is open and the gap is continuous along continuous
families (III-A)\@. On a frustration-free / LTQO stratum, the
Bravyi--Hastings--Michalakis and Michalakis--Zwolak stability theorems with the
Nachtergaele--Sims--Young bulk-gap result \cite{bhm-stability,michalakis-zwolak,nsy-bulkgap}
give a size-independent perturbation threshold $\varepsilon_0$: for
$\Phi$ in the stratum and a quasi-local direction $V$, the segment
$\set{\Phi+tV:0\le t<\varepsilon_0}\subseteq\Gap_{d,G}$ with gap bounded below by
$\gamma/2$. Hence $\Gap_{d,G}$ is condensed-open along the quasi-local direction class at
each stratum point.
\end{recollection}

\begin{recollection}[Uniform clustering; Part~III, Theorem~III-C \cite{paperGap}]
\label{rec:IIIc}
A uniformly gapped family over a profinite base, with a uniform Lieb--Robinson velocity
from \Cref{rec:Ic}, has exponential clustering with constants---hence a correlation
length---bounded uniformly over the base \cite{hastings-koma,nachtergaele-sims-clustering}.
\end{recollection}

\emph{Emergent structure.} Inverting $\Weq$ and stabilizing on $\Gap_{d,G}$ yields a
\emph{robust} $\pi_0$: a set of phases stable under quasi-local deformation, disorder
averaging, and finite-depth circuits. The robustness is exactly what
\Cref{rec:III,rec:IIIc} buy: openness so a phase label does not
change under small perturbations, clustering so it is detected locally. But
the openness holds only on the stratum where the stability hypotheses hold; Conjecture~III-1
asks for it on the full physical stratum, III-2 for descent along hulls, III-3 for the
identification of $\Weq$-components with Ogata's operator-algebraic phases
\cite{ogata-classification-review}.

\subsection[Arrow IV: stabilization and the phase spectrum]{Arrow IV: stabilization and the invertible condensed phase spectrum}
\label{sec:arrowIV}

The fourth arrow group-completes the invertible sector into a spectrum. On components
$\boxtimes$ makes $\Phases_{d,G}$ a commutative monoid with unit the trivial product
state; the short-range-entangled (SRE) phases are exactly its invertible elements.

\begin{recollection}[Group completion and the stabilization element; Part~IV \cite{paperEFT}]
\label{rec:groth}
For a commutative monoid $M$, the Grothendieck group $\Groth(M)$ has the universal
property that monoid maps $M\to A$ into abelian groups factor uniquely through
$M\to\Groth(M)$; $\gamma_M(a)=\gamma_M(b)$ iff $a+e=b+e$ for some $e\in M$, so
$\gamma_M$ is injective iff $M$ is cancellative. The element $e$ is the algebraic shadow
of ancilla stabilization.
\end{recollection}

\begin{recollection}[Kubota's $\Omega$-spectrum; Part~IV, after Kubota \cite{kubota-omega-spectrum,paperEFT}]
\label{rec:kubota}
There is an $\Omega$-spectrum $\mathit{IP}^\ast$, built from the operator-algebraic
formulation of invertible gapped quantum spin systems, whose homotopy groups are the
groups of invertible gapped systems in each dimension. It realizes Kitaev's proposal
\cite{kitaev-periodic-table} that invertible phases are the homotopy groups of a
spectrum and is the rigorous carrier for the homotopy of $\IPcond_{d,G}$.
\end{recollection}

Applying the recognition principle for grouplike $E_\infty$-spaces to the invertible
sector yields a connective spectrum, the invertible condensed phase spectrum
$\IPcond_{d,G}$ \cite{beaudry-etal}. Against it Part~IV sets the effective-field-theory
target: the Freed--Hopkins Anderson-dual bordism spectrum $\IFH$ classifying invertible
topological field theories \cite{freed-hopkins,freed-sre}.

\begin{recollection}[Hall conductance as a phase invariant; Part~IV, after Kapustin--Sopenko \cite{kapustin-sopenko-hall,kapustin-sopenko-noether,paperEFT}]
\label{rec:hall}
For a two-dimensional SRE lattice state the Hall conductance is locally computable, an
integer multiple of $e^2/h$, and constant on gapped phases, hence a homomorphism
$\Phases_{2,G}^\times\to\ZZ$; the higher Berry class generalizes it to families and
unifies it with the Thouless pump \cite{thouless-pump}.
\end{recollection}

\emph{Emergent structure.} Two things appear only at the spectrum level. First, the
\emph{higher homotopy}: $\pi_1$ of the phase object is loops of gapped Hamiltonians,
which implement adiabatic pumps and automorphisms of topological order
\cite{aasen-wang-hastings}; $\pi_n$ is $n$-parameter families and higher defects.
Second, the \emph{transition calculus} becomes a long-exact-sequence computation, the
relative charge $\rc$ appearing as a boundary map (\Cref{sec:transitions}). Part~IV's
Conjectures~IV-1 through IV-5 govern the condensed refinement, the lattice--EFT
comparison, the boundary-map identity, the compatibility of solidification with group
completion, and the renormalization functor.

\subsection{Arrow V: realizability}
\label{sec:arrowV}

The fifth arrow asks the converse question: which classes in $\IPcond_{d,G}$ are
produced by an explicit uniformly gapped lattice Hamiltonian? Part~V formalizes
realizability as essential surjectivity on $\pi_0$ of the comparison map $\real$ from
the stack of lattice models to the abstract spectrum, and separates realizability by
\emph{some} gapped local model from realizability by a local \emph{commuting-projector}
model.

\begin{recollection}[Cohomological realizability; Part~V, Theorem~V-A \cite{paperRealizability}]
\label{rec:VA}
For every finite $G$ and every $\omega\in H^{d+1}(G,\Uone)$ there is an explicit
$G$-symmetric local commuting-projector Hamiltonian $H_\omega$ with a unique
short-range-entangled ground state on any closed $d$-manifold, whose boundary carries the
anomalous $G$-action with obstruction $\omega$; $\omega\mapsto[H_\omega]$ is a
homomorphism $H^{d+1}(G,\Uone)\to\SRE_{d,G}$ injective on the group-cohomology subgroup
\cite{cglw-cohomology,else-nayak}.
\end{recollection}

\begin{recollection}[Low-dimensional completeness; Part~V, Theorem~V-B and Remark~V-B$'$ \cite{paperRealizability}]
\label{rec:VB}
For on-site finite symmetry in $d=1$, Ogata's $H^2(G,\Uone)$-valued index is a
\emph{complete} invariant of the symmetric phase
\cite{ogata-spt-chains,ogata-classification-review}, so $\real$ is a bijection onto
$H^2(G,\Uone)$ and realized $=$ invariant. In $d=2$ the situation is genuinely weaker:
Ogata's $H^3(G,\Uone)$-valued index \emph{exists} and is a well-defined invariant
\cite{ogata-h3-index}, and together with the group-cohomology models of \Cref{rec:VA} this
makes $\real$ \emph{surjective} onto the in-cohomology classes---so every such class is
realized---but whether the index is \emph{complete} (separates all $d=2$ SPT phases with
on-site finite symmetry) is open. We therefore do \emph{not} assert realized $=$ invariant
in $d=2$; only the surjectivity half is settled there.
\end{recollection}

\begin{recollection}[Commuting-projector no-gos; Part~V, Theorem~V-C \cite{kapustin-fidkowski,kapustin-spodyneiko-thermal,paperRealizability}]
\label{rec:VC}
A local commuting-projector Hamiltonian in two dimensions obeys two logically independent
no-gos. (i) \emph{Electric} (Kapustin--Fidkowski \cite{kapustin-fidkowski}): if it is
$U(1)$-symmetric, its zero-temperature electric Hall conductance vanishes, $\sigma_H=0$.
(ii) \emph{Thermal} (Kapustin--Spodyneiko \cite{kapustin-spodyneiko-thermal}): its chiral
central charge vanishes, $c_-=0$, with no symmetry hypothesis. Hence a $U(1)$-symmetric
phase with $\sigma_H\ne0$ (the integer quantum Hall states, Chern insulators) is excluded by
(i), and \emph{any} phase with $c_-\ne0$---including the bosonic $E_8$ state, which carries
no $U(1)$ charge (so $\sigma_H=0$) yet has $c_-=8$---is excluded by the chiral mechanism (ii),
not by $\sigma_H$.
\end{recollection}

\emph{Emergent structure.} The composition of all five arrows produces an
\emph{obstruction subgroup} $\Obs_{d,G}=\pi_0\IPcond_{d,G}/\im\real$ separating abstract
classes from lattice-realizable ones, with the chiral classes of \Cref{rec:VC} sitting
inside the commuting-projector obstruction. Part~V's Conjectures~V-1, V-2, V-3 describe
the image, the chiral classes, and profinite-family realizability.

\section{Modular composition}
\label{sec:composition}

The previous section walked the arrows in isolation. Here we make the \emph{composition}
explicit: how the modules compose, what emerges from composition that is invisible in any
single module, and what bounds the composition from outside.

\subsection{The composition diagram}
\label{sec:compdiagram}

Read as modules, the five Parts form a diagram in which each object is the input to the
next and each module contributes a functor. Write $\mathsf{Loc}$, $\mathsf{Obs}$,
$\mathsf{Gap}$, $\mathsf{Stab}$, $\mathsf{Real}$ for the five modules.

\begin{equation}
\label{eq:compdiagram}
\begin{tikzcd}[column sep=1.5em, row sep=2.6em]
\{\text{loc.\ int.}\}
  \arrow[r, "\mathsf{Loc}"]
& \Ham_{d,G}
  \arrow[r, "\mathsf{Gap}"]
  \arrow[d, "\mathsf{Obs}"']
& \Gap_{d,G}
  \arrow[r, "\mathsf{Stab}"]
& \Phase_{d,G}
  \arrow[r, "\boxtimes\text{, }\Groth"]
& \IPcond_{d,G}
  \arrow[d, "c_{d,G}"]
  \arrow[dl, dashed, "\real"'] \\
& (\qA,\ \cond{\St(\qA)},\ \Solid\Kalg)
&& \{\text{realized}\}
& \IFH
\end{tikzcd}
\end{equation}

Here $\mathsf{Loc}$, $\mathsf{Obs}$, $\mathsf{Gap}$, $\mathsf{Stab}$, and the
group-completion arrow are the functors of Parts~I, II, III, III--IV, and IV\@; the solid
arrows are functors the modules construct, with the caveat that $\mathsf{Stab}$ inverts
$\Weq$, whose generation from quasi-adiabatic continuation is Conjecture~III-3. The dashed
arrow $\real$ is the realization comparison of \Cref{sec:arrowV} (Part~V), whose essential
image is Conjecture~V-1; the vertical comparison $c_{d,G}:\IPcond_{d,G}\to\IFH$ to the
Freed--Hopkins target is Conjecture~IV-2. The observable data
$(\qA,\cond{\St(\qA)},\Solid\Kalg)$ produced by $\mathsf{Obs}$ feeds the homotopy of the
spectrum through the solid refinement of Conjecture~IV-1. The diagram commutes at the level
of $\pi_0$ exactly where those conjectures hold---in $d=1$ for on-site finite symmetry
(\Cref{rec:VB}), and for the free-fermion tenfold way
\cite{kitaev-periodic-table,altland-zirnbauer}---and its global commutativity is the Master
Conjecture of \Cref{sec:conjectures}.

\subsection{Emergent properties at each level}
\label{sec:emergent}

The characteristic feature of a modular composition is that structure emerges at each
level which is not a property of the parts. We collect the four emergences named in
\Cref{sec:spine}.

\begin{enumerate}
\item \textbf{Uniformity} (I$\circ$nothing). A single light cone, clustering length, and
  constant set over a whole profinite base: the precondition for every ``family'' below.
\item \textbf{Functorial disordered invariants} (II$\circ$I). Topological invariants of
  disordered families computed by solidification, matched to the physical configuration
  space rather than a smooth approximation of it.
\item \textbf{A robust phase set} (III$\circ$II$\circ$I). $\Phases_{d,G}=\pi_0\Shape\Phase_{d,G}$,
  stable under quasi-local deformation once $\Weq$ is inverted and ancillas added.
\item \textbf{Higher homotopy and an obstruction subgroup}
  (V$\circ$IV$\circ$III$\circ$II$\circ$I). Pumps and defects as $\pi_{\ge1}$, and
  $\Obs_{d,G}$ separating abstract from realizable.
\end{enumerate}

None of these emergences is forced by the module below it alone; each requires the
composite. That is the precise content of calling the program modular rather than
layered: the modules interoperate, and the interoperation is where the physics of
families, disorder, and defects lives.

\subsection{The two walls}
\label{sec:walls}

Two theorems bound the composition and must be respected at every level.

The \emph{undecidability wall} (\Cref{rec:cpw}) forbids any claim that $\Gap_{d,G}$ is
decidable or algorithmically presentable. It does not forbid stability results: openness
on a stratum, uniform clustering, and descent along a hull are all statements about the
\emph{neighborhood} of a gapped system, not tests for gappedness, and Part~III is careful
to scope them to strata with uniform constants. The wall is why the program is organized
around stability, not existence.

The \emph{no-go wall} (\Cref{rec:VC}), combining the electric Kapustin--Fidkowski and the
thermal Kapustin--Spodyneiko no-gos, forbids any claim that a chiral invertible class is
realized by commuting projectors. It bounds $\im\real^{\CPr}$ to the non-chiral sector and
makes the separation between ``exactly solvable'' and ``physically realizable''
structural, not incidental. Every realizability statement in the program excludes the
chiral commuting-projector case by hypothesis.

A third boundary is not a theorem but a scope limit: a spectrum classifies only
invertible order. General anyon theories and non-invertible phases need higher categories
of excitations (braided or modular tensor categories in $2{+}1$ dimensions and their
higher analogues), so a condensed treatment of them would require a condensed higher stack
of phases and defects, not a condensed $K$-theory spectrum. That extension is outside the
present program.

\section{The homotopy dictionary}
\label{sec:dictionary}

Classification by $\pi_0$ records only which phases exist. The higher homotopy of the
phase object carries more, and the program fixes a dictionary between homotopy-theoretic
and physical data. We quote it verbatim from the prospectus, now with each line anchored
to the module that supplies it.

\begin{equation}
\label{eq:dictionary}
\boxed{
\begin{aligned}
\pi_0 &= \text{phases},\\
\pi_1 &= \text{adiabatic pumps and phase automorphisms},\\
\pi_n &= \text{higher families and higher defects},\\
\Sig &= \text{gapless transition locus},\\
E^{q+1}(B,B\setminus\Sig) &= \text{relative charge of a transition}.
\end{aligned}}
\end{equation}

The first line is Arrow~III's emergent structure, made precise by Conjecture~III-3: the
components of $\Gap_{d,G}[\Weq^{-1}]$ are the operator-algebraic gapped phases of Ogata
\cite{ogata-classification-review}. The second and third are Arrow~IV's: loops and
higher families of gapped Hamiltonians realize pumps and defects, carried by the homotopy
of Kubota's spectrum \cite{kubota-omega-spectrum} and, conjecturally, by the higher
homotopy solid modules of $\IPcond_{d,G}$ (Conjecture~IV-1); the automorphism content of
$\pi_1$ is the subject of \cite{aasen-wang-hastings}, whose $\pi_1,\pi_2,\pi_3$
computations are themselves conjectural. The last two lines are the transition calculus,
developed next.

\section{Phase transitions and the relative-charge formalism}
\label{sec:transitions}

This section develops the fourth and fifth lines of \eqref{eq:dictionary}: what a
transition is, and what its charge is. The development is honest obstruction theory (the
long exact sequence of a pair and excision, valid for any generalized cohomology), and
the only conjectural step is the identification of the abstract boundary map with the
spectral boundary map of $\IPcond_{d,G}$, which is Conjecture~IV-3 lifted to the program.

\subsection{The gapped locus, the discriminant, and the phase label}
\label{sec:locus}

Let $B$ be a parameter space (couplings, fields, pressures, disorder configurations,
boundary conditions), regarded as a condensed object $\cond{B}$, and let
$f:\cond{B}\to\Ham_{d,G}$ be a family of systems.

\begin{definition}[gapped locus and discriminant]
\label{def:locus}
The \emph{gapped locus} of $f$ is the pullback
\[
U_f=\cond{B}\times_{\Ham_{d,G}}\Gap_{d,G},
\]
and the \emph{gapless (critical) locus} or \emph{discriminant} is its complement
$\Sig_f=\cond{B}\setminus U_f$. On $U_f$ the phase label
\[
\nu_f:U_f\longrightarrow\pi_0(\Phase_{d,G})=\Phases_{d,G}
\]
is locally constant.
\end{definition}

That $\nu_f$ is locally constant on $U_f$ is exactly \Cref{rec:III}: the gapped locus is
condensed-open on the physical stratum, and openness of each phase's preimage is what
``locally constant'' means. A topological phase transition is a path in $B$ whose
endpoints carry different values of $\nu_f$; such a path cannot stay in $U_f$ and must
meet $\Sig_f$. This is the boxed slogan \eqref{eq:slogan2}.

\subsection{The relative charge as an obstruction}
\label{sec:relcharge}

Suppose the invariant is valued in a (condensed) generalized cohomology theory $E$; in
the invertible case $E$ is represented by $\IPcond_{d,G}$. On the gapped locus the
invariant is a class $\nu\in E^q(U_f)$. The question of the transition is whether $\nu$
extends over the critical set to all of $B$. Obstruction theory answers it through the
long exact sequence of the pair $(B,U_f)$:
\begin{equation}
\label{eq:les}
\cdots\to E^q(B)\xrightarrow{\ j^\ast\ }E^q(U_f)\xrightarrow{\ \partial\ }
E^{q+1}(B,U_f)\xrightarrow{\ \ }E^{q+1}(B)\to\cdots,
\end{equation}
with $j:U_f\hookrightarrow B$ the inclusion.

\begin{definition}[relative transition charge]
\label{def:relcharge}
The \emph{relative transition charge} of the family $f$ carrying the invariant $\nu\in
E^q(U_f)$ is
\[
\rc:=\partial\nu\in E^{q+1}(B,U_f)=E^{q+1}(B,B\setminus\Sig_f),
\]
the image of $\nu$ under the connecting map of \eqref{eq:les}.
\end{definition}

\begin{proposition}[extension obstruction]
\label{prop:obstruction}
$\rc=0$ if and only if $\nu$ is the restriction of a class on all of $B$; that is,
$\rc$ is the obstruction to extending the phase invariant across the discriminant. In
particular, if $\rc\ne0$ the family has a genuine transition: no global invariant
restricts to $\nu$ on $U_f$.
\end{proposition}

\begin{proof}
Exactness of \eqref{eq:les} at $E^q(U_f)$: $\nu\in\im(j^\ast)$ if and only if
$\partial\nu=0$. A class in $\im(j^\ast)$ is by definition the restriction of a class on
$B$. This is the standard long exact sequence of a pair in a generalized cohomology
theory and requires nothing beyond the Eilenberg--Steenrod axioms that $E$ satisfies as
a spectrum-represented theory.
\end{proof}

\subsection{Locality: linking spheres}
\label{sec:linking}

The relative group $E^{q+1}(B,B\setminus\Sig_f)$ is local along $\Sig_f$. When $\Sig_f$
is a reasonable (say, closed, locally-flat, finite-codimension) condensed subobject with
components $\{\Sig_f^{(\alpha)}\}$, excision identifies the relative cohomology with a sum
of contributions supported near the components:
\begin{equation}
\label{eq:excision}
E^{q+1}(B,B\setminus\Sig_f)\ \cong\ \bigoplus_\alpha E^{q+1}\bigl(N_\alpha,N_\alpha\setminus\Sig_f^{(\alpha)}\bigr),
\end{equation}
$N_\alpha$ a tubular neighborhood of $\Sig_f^{(\alpha)}$. If $\Sig_f^{(\alpha)}$ has
codimension $c$ with an $E$-oriented normal bundle, the Thom isomorphism identifies each
summand with a degree-shifted absolute group of the component,
\begin{equation}
\label{eq:linking}
E^{q+1}\bigl(N_\alpha,N_\alpha\setminus\Sig_f^{(\alpha)}\bigr)\ \cong\
E^{q+1-c}\bigl(\Sig_f^{(\alpha)}\bigr),
\end{equation}
the Thom class shifting degree by the codimension $c$. The local charge is read off on a
linking sphere: over a point $x\in\Sig_f^{(\alpha)}$ the unit sphere $S^{c-1}_{\mathrm{link}}$
of the normal fibre lies in $B\setminus\Sig_f$ and links the component, and the restriction
$\nu|_{S^{c-1}_{\mathrm{link}}}\in E^{q}\bigl(S^{c-1}_{\mathrm{link}}\bigr)$ has reduced part
$\widetilde E^{\,q}(S^{c-1})\cong E^{q-c+1}(\mathrm{pt})$ equal to that local charge: the
same degree shift as \eqref{eq:linking}. In words: the relative charge decomposes as a sum
of local charges, one per component of the critical set, each detected by evaluating $E$ on
a small sphere linking the component. This is the precise form of the physical picture that
a band degeneracy (a Dirac node, a Weyl point) acts as a source or sink of topological
charge, now valid for a generalized cohomology theory and an interacting phase spectrum
rather than only for Chern numbers of free bands. This relative, defect-localized reading
has a free-fermion precedent: Teo and Kane classify topological defects by the $K$-theory of
a sphere linking the defect \cite{teo-kane}, within the $K$-theoretic classification of
topological phases \cite{thiang-ktheory}; the condensed formulation of
\Cref{def:relcharge} is the extension of that relative-charge picture to a generalized
theory represented by $\IPcond_{d,G}$. The low-codimension case $c=1$ (a wall
$\Sig_f$ separating two gapped regions) reduces the linking sphere to $S^0$: two points,
one in each phase, and the local charge is the \emph{difference} of the two phase
labels: the jump of $\nu_f$ across the wall.

\subsection{The spectral boundary map}
\label{sec:boundary}

\Cref{prop:obstruction} and \eqref{eq:linking} are unconditional facts about any
$E$. The program's substantive claim is that, for $E=\IPcond_{d,G}$-cohomology, the
connecting map $\partial$ of \eqref{eq:les} is the boundary map of the cofiber sequence
of the pair for the spectrum $\IPcond_{d,G}$, so the transition charge is computed by a
single long exact sequence in every dimension and every generalized theory
simultaneously. This is Conjecture~IV-3, restated at the program level as
Conjecture~VI-3 in \Cref{sec:conjectures}. Its low-dimensional shadow (the SSH winding
jump, the Dirac node as a unit source of charge) is classical; the conjecture is that
this is systematically the spectral boundary map.

\section{The SSH transition end to end}
\label{sec:ssh}

We now run one example through all five modules, to show that the composition is not
merely formal. The Su--Schrieffer--Heeger chain \cite{ssh-1979} is the smallest system in
which every module has something to say, and its winding number is a degree-$\ge0$ class,
so it sits safely on the nonnegative side of the Bott-inversion caveat of \Cref{rec:aoki}.

\begin{example}[SSH through the spine]
\label{ex:ssh}
The SSH Bloch Hamiltonian and its off-diagonal function are
\[
H(k;t_1,t_2)=(t_1+t_2\cos k)\,\sigma_x+(t_2\sin k)\,\sigma_y,
\qquad q(k)=t_1+t_2 e^{ik},
\]
with spectrum $E_\pm(k)=\pm\abs{q(k)}$.
\end{example}

\begin{proposition}[gap and winding of the SSH chain]
\label{prop:sshgap}
The chain is gapped if and only if $\abs{t_1}\ne\abs{t_2}$; the gap closes on the
discriminant $\Sig=\set{(t_1,t_2):\abs{t_1}=\abs{t_2}}$. On the gapped locus the winding
number of $q:S^1\to\CC^\times$ is
\[
\nu=\begin{cases}1,&\abs{t_1}<\abs{t_2},\\ 0,&\abs{t_1}>\abs{t_2},\end{cases}
\]
up to orientation, and jumps by $\pm1$ across $\Sig$.
\end{proposition}

\begin{proof}
$E_\pm(k)=\pm\abs{q(k)}$ vanishes for some $k$ iff $q(k)=0$, i.e.\ iff $t_1+t_2e^{ik}=0$
for some $k$, which happens iff $\abs{t_1}=\abs{t_2}$; hence the gap is open exactly off
$\Sig$. On the gapped locus $q$ misses the origin, so its winding number is defined; for
$\abs{t_1}>\abs{t_2}$ the curve $k\mapsto t_1+t_2e^{ik}$ is a circle of radius
$\abs{t_2}$ about $t_1\ne0$ not enclosing the origin (winding $0$), and for
$\abs{t_1}<\abs{t_2}$ it encloses the origin once (winding $1$). The two regions are
separated by $\Sig$, across which the winding jumps by $\pm1$.
\end{proof}

Now the five modules, in order.

\emph{Part~I (dynamics).} The SSH couplings define a finite-range interaction of finite
$\Ffun$-norm, so the family lives in $\cond{\BF}$ and, by \Cref{rec:Ic}, over any compact
region of the $(t_1,t_2)$-plane away from $\Sig$ carries one Lieb--Robinson velocity. The
parameter torus and Brillouin circle are replaced by their condensations
$\cond{B},\cond{S^1}$; the light cone is uniform over the probe.

\emph{Part~II (algebra and $K$-theory).} SSH is a class $\mathit{AIII}$ system in $d=1$,
whose tenfold-way entry is $\ZZ$ \cite{kitaev-periodic-table,altland-zirnbauer}, the
winding number valued in $K^1$ of the observable algebra. Because the winding is a
nonnegative-degree class, \Cref{rec:aoki} applies without invoking the Bott inversion:
the solid invariant $\Solid(\Kalg(\cond{A}))$ computes it directly, and its provenance is
condensed even though the number is the classical winding.

\emph{Part~III (gap).} \Cref{prop:sshgap} is exactly the finite-volume/thermodynamic gap
statement of \Cref{rec:III} for this family: the gapped locus is
$U=\set{\abs{t_1}\ne\abs{t_2}}$, open, and the phase label $\nu$ is locally constant on
it, with the discriminant $\Sig=\set{\abs{t_1}=\abs{t_2}}$.

\emph{Part~IV (stabilization and EFT).} The low-energy theory at the gap-closing
$\abs{t_1}=\abs{t_2}$, $k=\pi$ is a massive $1{+}1$-dimensional Dirac fermion, and the
winding number is the sign of the Dirac mass---the deformation class of the Dirac
effective field theory. So the comparison map $c_{1,\mathit{AIII}}$ sends the SSH phase to
its EFT class and the two agree: Conjecture~IV-2 holds at the level of $\pi_0$ in this
box.

\emph{Part~V (realizability).} The nontrivial SSH phase is short-range-entangled and
realized by the explicit dimerized chain; being non-chiral ($\sigma_H=0$ and $c_-=0$), it is
excluded by neither no-go of \Cref{rec:VC} and is in fact commuting-projector realizable in
the flat-band (fully dimerized) limit. It sits in $\im\real$.

\emph{The transition charge.} In the parameter plane $B=(t_1,t_2)$ the invariant relevant to
the transition calculus is not the momentum-space $K^1$ class of the tenfold way but the
\emph{winding number} it produces: a locally constant $\ZZ$-valued label on the gapped locus,
an element $\nu\in K^0(U)$ of degree $q=0$. The discriminant $\Sig=\set{\abs{t_1}=\abs{t_2}}$
is the pair of lines $t_1=\pm t_2$; away from their intersection it has codimension $c=1$, so
the linking sphere of \eqref{eq:linking} is $S^0$ and the local charge is read in the reduced
group $\widetilde K^0(S^0)\cong\ZZ$---the winding difference $\pm1$ across the wall, in
agreement with the degree shift $E^{q-c+1}(\mathrm{pt})=K^0(\mathrm{pt})=\ZZ$ of
\eqref{eq:linking}. (At the origin $(0,0)$ the two lines cross, the gap closes at $k=0$ and
$k=\pi$ simultaneously, and the codimension jumps to $2$; the generic crossing is the
codimension-one one just described.) By \Cref{prop:obstruction} the relative charge
$\rc\in K^{1}(B,U)$ is nonzero precisely because the winding label does not extend
across $\Sig$: the transition is real, and its charge is the unit jump. Conjecture~VI-3 is the
assertion that this $\partial$ is the spectral boundary map of $\IPcond_{1,\mathit{AIII}}$;
here it is the elementary winding jump, and the two agree.

\section{Disorder and the profinite thread}
\label{sec:disorder}

The profinite disorder thread is where condensed mathematics does work that ordinary
smooth-parameter topology cannot, and it runs through all five modules. We collect it.

\subsection{The hull as a profinite probe}

If $Q$ is a finite set of local configurations, the configuration space
\begin{equation}
\label{eq:hull}
\dis=Q^{\ZZ^d}
\end{equation}
is compact, totally disconnected, hence profinite; this is the disorder (or tiling) hull of
Bellissard's noncommutative geometry of aperiodic media \cite{bellissard-ncg-qhe,kellendonk-tilings}.
(We follow Part~II in writing the
alphabet $Q$ rather than the prospectus's $\Ffun$, to avoid collision with the
Nachtergaele--Sims--Young $\Ffun$-function of Part~I; the two never denote the same
object.) A disordered Hamiltonian family $\omega\mapsto H_\omega$ is literally an
$\dis$-point of the moduli stack, $H\in\Ham_{d,G}(\dis)$, and finite quotients of $\dis$
are finite-resolution disorder data. The sheaf condition of \eqref{eq:condensation} is
the compatibility of families and invariants across finite approximations, so the
profinite probes match the physical configuration space rather than repackage a smooth
one.

\subsection{The thread through the modules}

\begin{itemize}
\item \textbf{Part~I} (Conjecture~I-3): the Lieb--Robinson estimates over $\dis$ should be
  the right Kan extension of their restrictions to finite quotients: quasi-local dynamics
  determined by finite-resolution data. Its unconditional shadow is the uniform light cone
  of \Cref{rec:Ic}.
\item \textbf{Part~II} (\Cref{rec:crossed}, Conjecture~II-2): the observable algebra is the
  crossed product $C(\dis)\rtimes_T\ZZ^d$, its condensation functorial in finite quotients,
  and the solid invariant should agree with operator $K$-theory naturally in $\dis$.
\item \textbf{Part~III} (Conjecture~III-2): uniform gappedness over $\dis$ should be a
  closed condition on the inverse system: the uniformly gapped families the inverse limit
  of the finite-resolution ones at a fixed $\gapbound$.
\item \textbf{Part~IV}: parametrized invariants over $\dis$, the base level (L2) of the
  lattice--EFT comparison Conjecture~IV-2.
\item \textbf{Part~V} (Conjecture~V-3): realizability over $\dis$ is disorder-robust iff a
  condensed-cohomological descent obstruction $o(c)\in H^1_{\mathrm{cond}}(\dis,\cg{\Aut}(c))$
  vanishes, detected by finite quotients.
\end{itemize}

The common shape of all five is descent along the finite quotients of $\dis$. That the
five descent statements are facets of one principle is Conjecture~VI-4.

\section{The program-level conjectures}
\label{sec:conjectures}

We now state the open claims of the program. Each is built as the coherent join of
module-level conjectures from Parts~I--V, and none is proved. We write them as named
Conjectures~VI-1 through VI-4; the first is the Master Conjecture.

\begin{progconj}[Master Conjecture: the modules glue into one condensed higher stack]
\label{conj:VI1}
The five module-level structural conjectures hold simultaneously and coherently:
$\Ham_{d,G}$ is a condensed higher stack (Conjecture~I-1); positivity and the
$C^\ast$-identity cut it out as a closed condensed substack of a formal-interaction stack
(Conjecture~II-1); $\Gap_{d,G}$ is an open condensed substack on the physical stratum with
the uniform-gap sheaf condition detected by finite quotients (Conjecture~III-1); the
stabilized invertible sector is a connective condensed/solid spectrum $\IPcond_{d,G}$
whose shape is a connective cover of Kubota's $\mathit{IP}^\ast$ (Conjecture~IV-1); and
the realizability image is the short-range-entangled subspectrum (Conjecture~V-1).
Consequently the boxed sequence \eqref{eq:boxseq} is a diagram of condensed higher stacks
and spectra, $\Phases_{d,G}=\pi_0\Shape(\Phase_{d,G})$ is the operator-algebraic set of
gapped phases, and $\IPcond_{d,G}$ is the invertible phase spectrum of the whole.
\end{progconj}

\begin{progconj}[microscopic $=$ field-theoretic]
\label{conj:VI2}
Under short-range-entanglement hypotheses the comparison map
$c_{d,G}:\IPcond_{d,G}\to\IFH$ to the Freed--Hopkins invertible-TQFT spectrum is an
equivalence after solidification/completion (Conjecture~IV-2), and its $\pi_0$ image
coincides with the realized SRE subspectrum $\im\real=\SRE_{d,G}$ (Conjecture~V-1). Thus
the microscopic classification of $\Phases_{d,G}^\times$ and the bordism classification of
$\pi_0\IFH$ agree exactly on the realizable classes, with the chiral part realizable by
non-commuting gapped models but never by commuting projectors (Conjecture~V-2,
\Cref{rec:VC}). The statement holds at $\pi_0$ in $d=1$ for on-site finite symmetry
(\Cref{rec:VB}) and for the free-fermion tenfold way, and is open in general.
\end{progconj}

\begin{progconj}[the relative charge is a spectral boundary map]
\label{conj:VI3}
For $E=\IPcond_{d,G}$-cohomology, the connecting map $\partial$ of \eqref{eq:les} is the
boundary map of the cofiber sequence of the pair $(B,U_f)$ for the spectrum
$\IPcond_{d,G}$ (Conjecture~IV-3). Hence the relative transition charge
$\rc=\partial\nu\in E^{q+1}(B,B\setminus\Sig_f)$ of \Cref{def:relcharge} is computed by a
single long exact sequence in every dimension and generalized theory, and its
linking-sphere localization \eqref{eq:linking} is the systematic form of ``a band
degeneracy is a source of topological charge.'' The SSH box (\Cref{ex:ssh}) is the
codimension-one, winding-number instance.
\end{progconj}

\begin{progconj}[all program invariants descend along disorder hulls]
\label{conj:VI4}
For a profinite disorder hull $\dis=Q^{\ZZ^d}=\varprojlim_i\dis_i$, the disorder facets
of the five modules---Lieb--Robinson descent (I-3), naturality of the solid invariant
(II-2), uniform-gap descent (III-2), parametrized comparison over $\dis$ (IV-2 at level
L2), and profinite-family realizability (V-3)---are facets of one descent principle: every
program invariant over $\dis$ is the right Kan extension of its restrictions to the finite
quotients $\dis_i$, and uniform gappedness, the solid invariant, and realizability are all
detected by finite quotients with a common bound. This is the precise sense in which
``profinite probes match the configuration space.''
\end{progconj}

\subsection{Consolidated open problems}
\label{sec:opentable}

\Cref{tab:conjectures} lists the nineteen module-level conjectures of Parts~I--V, with the
module each depends on and the program-level conjecture it feeds. The dependency column
records which \emph{earlier} module's output a conjecture presupposes; the ``feeds''
column records which Conjecture~VI-$k$ consolidates it. The table is the program's
to-do list.

\begin{table}[htbp]
\centering
\small
\renewcommand{\arraystretch}{0.95}
\caption{The nineteen module-level conjectures of Parts~I--V and their dependencies.
``Feeds'' names the program-level Conjecture~VI-$k$ that consolidates each.}
\label{tab:conjectures}
\begin{tabular}{@{}l >{\raggedright\arraybackslash}p{7.7cm} l l@{}}
\toprule
ID & Statement (abbreviated) & Depends on & Feeds \\
\midrule
I-1 & $\Ham_{d,G}$ is a condensed higher stack (descent of the automorphism groupoid) & --- & VI-1 \\
I-2 & Quasi-local automorphisms form a condensed group; QAC is internal path-lifting & I-1 & VI-1 \\
I-3 & Lieb--Robinson estimates descend along a disorder hull & I-1 & VI-4 \\
\midrule
II-1 & Positivity $+$ $C^\ast$-identity cut $\Ham_{d,G}$ out as a closed condensed substack & I-1 & VI-1 \\
II-2 & Solid invariant $=$ operator $K$ of the crossed product, naturally in $\dis$ & II-1 & VI-4 \\
II-3 & Real/$KKO$ refinement recovering the $KO$-graded periodic table & II-1 & VI-2 \\
II-4 & Solid state space is a classifying object for condensed representations & II-1 & VI-1 \\
II-5 & Condensed enhancement of the split property / DHR superselection & II-1 & VI-1 \\
\midrule
III-1 & $\Gap_{d,G}$ open on the physical stratum; sheaf condition detected by finite quotients & II-1 & VI-1 \\
III-2 & Uniform-gap descent along profinite hulls (closed on the inverse system) & III-1 & VI-4 \\
III-3 & QAC path components realize $\Weq$; $\pi_0\Shape(\Gap[\Weq^{-1}])=$ Ogata phases & III-1 & VI-1 \\
\midrule
IV-1 & $\IPcond_{d,G}$ is a connective condensed/solid spectrum covering Kubota's $\mathit{IP}^\ast$ & III-3 & VI-1 \\
IV-2 & Lattice--EFT comparison $c_{d,G}:\IPcond_{d,G}\to\IFH$ is an equivalence after completion & IV-1 & VI-2 \\
IV-3 & Relative charge $\rc$ is the spectral boundary map of $\IPcond_{d,G}$ & IV-1 & VI-3 \\
IV-4 & Solidification commutes with group completion & IV-1 & VI-1 \\
IV-5 & Renormalization is a filtered pro-endofunctor computing the EFT on the SRE stratum & IV-1 & VI-2 \\
\midrule
V-1 & Realizability image $=$ SRE subspectrum; $\Obs_{d,G}=\pi_0\IPcond_{d,G}/\im\real$ & IV-2 & VI-1 \\
V-2 & Chiral classes realizable by non-commuting gapped, never commuting-projector & V-1 & VI-2 \\
V-3 & Profinite-family realizability iff a condensed descent obstruction vanishes & V-1 & VI-4 \\
\bottomrule
\end{tabular}
\end{table}

Two features of \Cref{tab:conjectures} are worth naming. First, the dependency graph is a
tree rooted at I-1: every conjecture presupposes the stack structure of the ground floor,
which is why Conjecture~I-1 is the true foundation and the Master Conjecture VI-1 leads
with it. Second, the ``feeds'' column shows the program-level conjectures are not
independent wishes but bundles: VI-1 gathers the ten structural conjectures, VI-2 the
four comparison/refinement conjectures, VI-3 the single boundary-map conjecture, and VI-4
the four descent conjectures ($10+4+1+4=19$). Proving any program-level conjecture means proving its whole
bundle.

\FloatBarrier
\section{Contributions and limitations}
\label{sec:contributions}

\subsection{Contributions}

The framework contributes three structural advantages, none of which is a new numerical
invariant. First, one category holds ordinary continuous families, profinite disorder
families, compact inverse limits, and topological symmetry groups at once; the passage
\eqref{eq:condensation} is fully faithful, so nothing classical is lost. Second, condensed
abelian groups and solid modules supply exact derived operations and descent where
ordinary topological groups and modules are poorly behaved: the environment
\Cref{rec:aoki} needs to make operator $K$-theory a condensed invariant. Third, the sheaf
and stack viewpoint organizes local families, defects, interfaces, boundary conditions,
symmetry actions, and gluing uniformly, which is what the transition calculus of
\Cref{sec:transitions} and the disorder thread of \Cref{sec:disorder} exploit.

On novelty we are precise, because the reviewer should be. Condensed and pyknotic
mathematics \cite{clausen-scholze-condensed,barwick-haine-pyknotic} and the solidification
bridge \cite{aoki-solidification} are established in pure mathematics; the
moduli/space-of-states viewpoint on gapped systems is established in physics
\cite{hsin-wang-moduli,beaudry-etal,kubota-omega-spectrum,denittis-states,denittis-rendel-weyl}.
What is new is the \emph{synthesis}: no prior work applies condensed or pyknotic machinery
to topological phases, and the individual analytic ingredients (Lieb--Robinson bounds,
$C^\ast$-positivity, gap stability, stabilization, realizability) are not new in
themselves. The value is unificatory and organizational, and the honest content is the
Master Conjecture and its bundle.

\subsection{What it does not do}

By itself the paradigm does not solve the hard analysis of many-body theory: one still
needs locality estimates (Part~I), positivity and $C^\ast$-norm conditions (Part~II), the
existence and stability of a thermodynamic gap (Part~III), the lattice--EFT correspondence
(Part~IV), and physical realizability (Part~V), and each is a genuine theorem or a genuine
conjecture, not a corollary of the formalism. It does not decide the gap: the undecidability
wall (\Cref{rec:cpw}) is permanent. It does not realize chiral phases by commuting
projectors: the no-go wall (\Cref{rec:VC}) is permanent. It does not, by construction,
classify non-invertible topological order, which needs a condensed higher stack of phases
and defects rather than a spectrum. And it does not produce a different Chern number: the
winding of \Cref{prop:sshgap} is the classical winding; only its provenance is condensed.
The genuinely new thesis is the single environment, and the genuinely open mathematics is
Conjectures~VI-1 through VI-4.

\section{Conclusion}
\label{sec:conclusion}

The six papers of this series are one program with five modules and a spectrum-level
payoff. Part~I makes the interaction space condensed and the dynamics a base-uniform
morphism; Part~II attaches observable algebras, a compact condensed state space, and the
solid $K$-invariant with its Bott-inversion bookkeeping; Part~III cuts out the uniformly
gapped substack and stabilizes it against the undecidability wall; Part~IV group-completes
the invertible sector into $\IPcond_{d,G}$ and confronts the field-theoretic
classification; Part~V asks which classes a lattice realizes, bounded by the
commuting-projector no-gos of Kapustin--Fidkowski and Kapustin--Spodyneiko. Composed, the
modules give the boxed sequence
\eqref{eq:boxseq}, the homotopy dictionary \eqref{eq:dictionary}, and the transition
calculus in which a transition is the crossing of the discriminant and its charge is the
relative class $\rc$ localized by linking spheres.

We have been careful to keep the ledger honest. The module results are theorems, cited to
Parts~I--V\@. The program results (that the modules glue into one condensed higher stack
with $\IPcond_{d,G}$ its invertible spectrum, that microscopic and field-theoretic
classifications agree, that the relative charge is a spectral boundary map, that all
invariants descend along disorder hulls) are Conjectures~VI-1 through VI-4, and
\Cref{tab:conjectures} records exactly what each would need. The program is not a completed
classification and does not claim to be. It is a single categorical and homological
environment in which continuous families, profinite disorder, analytic completions,
operator $K$-theory, symmetry, stacking, defects, and phase-transition loci can be treated
at once, and a precise list of what remains to prove.

\section*{Code availability}
{\hbadness=10000
\begin{sloppypar}
The formal-verification suites accompanying the six-paper series are at
\href{https://github.com/YonedaAI/topological-phases-of-matter}{github.com/\allowbreak YonedaAI/\allowbreak topological-phases-of-matter};
each module's code lives in a per-topic subdirectory (for example
\texttt{src/\allowbreak bordism-realizability/} and
\texttt{src/\allowbreak lattice-eft-\allowbreak equivalence/}). This synthesis introduces no
new code of its own; it composes the results verified in Parts~I--V.
\end{sloppypar}}

\section*{Acknowledgements}
This is Part~VI of a six-part series by the author and The YonedaAI Collaboration; it
depends entirely on the analytic and topological content of Parts~I--V and on the
foundational work cited throughout.

\bibliographystyle{plain}
\bibliography{references}

\end{document}
